\chapter{Opis struktury projektu}
\label{cha:Opis struktury projektu}

\section{Technologie wykorzystane w projekcie}

Projekt został zaimplementowany w Pythonie w~wersji 3.12 z wykorzystaniem biblioteki pygame 2.6.1. 
Środowiskiem programistycznym był PyCharm Community Edition. Repozytorium kodu jest organizowane w jeden katalog 
roboczy.
 
Architektura gry opiera się na pętli zdarzeń pygame z zachowaniem stałej liczby klatek (60 FPS). Każda klatka obejmuje:
\begin{enumerate}
    \item Obsługę zdarzeń klawiatury (input gracza)
    \item Aktualizację AI przeciwników
    \item Aktualizację pozycji pocisków i~detekcję kolizji
    \item Renderowanie sceny na ekranie
\end{enumerate}

\section{Struktura plików źródłowych}

Projekt składa się z następujących plików źródłowych:
 
\begin{itemize}
    \item \texttt{main.py} - punkt wejścia aplikacji; obsługa menu, wybór trybu, pętla gry, 
    zarządzanie stanami rozgrywki, spawnowanie power-upów i przeciwników
    \item \texttt{tank.py} - klasa bazowa Tank reprezentująca czołg (gracza lub przeciwnika); 
    obsługa ruchu, rotacji, strzelania, kolizji, power-upów, systemu zdrowia
    \item \texttt{enemy.py} - klasa EnemyTank dziedzicząca po Tank; implementacja AI opartej 
    na Minimax z przycinaniem alfa-beta dla trybu Deathmatch
    \item \texttt{enemy2.py} - klasa EnemyTankCP dziedzicząca po Tank; implementacja AI opartej 
    na A* dla trybu Capture Point
    \item \texttt{bullet.py} - klasa Bullet reprezentująca pocisk; obsługa ruchu, animacji, detekcji wyjścia poza ekran
    \item \texttt{walls.py} - klasy ścian zniszczalnych Wall oraz niezniszczalnych IndestructibleWall
    \item \texttt{powerup.py} - klasa PowerUp reprezentująca bonus na mapie (życie, pancerz, silny strzał)
    \item \texttt{map\_loader.py} - funkcje wczytujące mapy z plików tekstowych .txt z folderów Maps (Deathmatch) i {Maps2} (Capture Point)
\end{itemize}
 
Mapa jest zapisana w pliku tekstowym każdy znak reprezentuje 
jeden kafelek 40x40 pikseli, \# to ściana niezniszczalna, X to ściana zniszczalna, spacja to wolne pole. 
Pozwala to na łatwe dodawanie nowych map bez modyfikacji kodu.

\section{Tryb Deathmatch - algorytm Minimax z przycinaniem alfa-beta}
\subsection{Idea algorytmu Minimax}

Minimax jest klasycznym algorytmem decyzyjnym przeznaczonym dla gier dwuosobowych z pełną informacją i ruchem naprzemiennym. 
Algorytm zakłada, że obaj gracze grają optymalnie: jeden (MAX) dąży do maksymalizacji oceny pozycji, drugi (MIN) do jej minimalizacji.
 
Drzewo Minimax o głębokości d reprezentuje sekwencję d kolejnych ruchów. Korzeń drzewa odpowiada 
bieżącej pozycji, każdy poziom poniżej reprezentuje jeden hipotetyczny ruch, na przemian bota (warstwa MAX) i gracza (warstwa MIN). 
W liściach drzewa (głębokość 0) wywoływana jest funkcja oceny pozycji $\text{eval}(s)$, której wartość propaguje się w górę 
drzewa zgodnie z regułą:
 
\begin{equation}
\text{Minimax}(s) =
\begin{cases}
\text{eval}(s) & \text{gdy } d = 0 \\
\max_{s' \in \text{succ}(s)} \text{Minimax}(s') & \text{gdy gracz to MAX} \\
\min_{s' \in \text{succ}(s)} \text{Minimax}(s') & \text{gdy gracz to MIN}
\end{cases}
\end{equation}
 
W zaimplementowanym projekcie głębokość drzewa wynosi d=4, czyli bot przewiduje dwa swoje ruchy i dwa odpowiedzi gracza. 
Strategia jest z natury defensywna - bot wybiera ruch, który jest dla niego najlepszy nawet w najgorszym scenariuszu, co eliminuje
błędy typu 'wskoczę na linię ognia, bo akurat gracz patrzy w drugą stronę'.
 
\subsection{Funkcja oceny pozycji}

Funkcja oceny \_eval składa się z czterech komponentów połączonych liniowo:
 
\begin{equation}
f(b, p, \text{cost}, \text{pu}) = -2 \cdot d_{Ch}(b, p) + B_{LOS} + B_{PU} - 2 \cdot \text{cost}_\text{acc}
\end{equation}
 
gdzie:
 
\textbf{Odległość Czebyszewa do gracza} z wagą $-2$:
\begin{equation}
d_{Ch}(a, b) = \max(|a_x - b_x|, |a_y - b_y|)
\end{equation}
 
Dystans Czebyszewa jest lepszą miarą niż odległość Manhattan dla ruchu w siatce 4-kierunkowej, 
ponieważ uwzględnia naturalne obwiednie ruchów (czołg może poruszać się w osi X i Y niezależnie). 
Im bliżej gracza, tym wyższa ocena, co motywuje bota do aktywnego ścigania przeciwnika.
 
\textbf{Bonus linii strzału (LOS)} $B_{LOS} = +80$ przyznawany, gdy bot i gracz znajdują się w tej samej osi 
(z tolerancją 16 pikseli) oraz między nimi nie ma niezniszczalnej ściany blokującej tor pocisku. 
Jest to bardzo silny sygnał: pozycja umożliwiająca czysty strzał jest mocno premiowana.
 
\textbf{Bonus power-upa} $B_{PU} \in [0, 40]$ przyznawany, gdy gracz jest daleko ($d_{Ch}(b, p) > 120$~px) 
a najbliższy power-up blisko ($d_{Ch}(b, \text{pu}) < 100$~px). Wartość bonusu skaluje się odwrotnie z 
odległością: $B_{PU} = \max(0, 40 - d_{Ch}(b, \text{pu})/3)$. Pozwala botowi zbierać użyteczne bonusy w momentach, gdy 
gonienie gracza nie ma sensu.
 
\textbf{Kara za koszt drogi} $-2 \cdot \text{cost}_\text{acc}$, gdzie $\text{cost}_\text{acc}$ to suma kosztów przejść 
przez zniszczalne ściany na symulowanej trasie. Każda zniszczalna ściana to koszt $\text{COST\_DEST} = 8$. 
Komponent ten zniechęca bota do planowania ruchów wymagających niszczenia ścian, jeśli istnieje alternatywna trasa.
 
Dodatkowo, w turze gracza w symulacji uwzględniony jest szczegół: jeśli na danej hipotetycznej 
pozycji gracz miałby linię strzału do bota, do oceny tej pozycji dodawana jest kara $-60$. 
Symulacja 'gracz strzela zamiast uciekać' wymusza na Minimaxie aktywne unikanie pozycji wystawiających bota 
na ostrzał, nie tylko maksymalizację bliskości.

\subsection{Przycinanie alfa-beta}

Przycinanie alfa-beta jest optymalizacją algorytmu Minimax pozwalającą pominąć całe gałęzie drzewa, 
które nie wpłyną na ostateczną decyzję. Algorytm utrzymuje dwie wartości graniczne:
 
\begin{itemize}
    \item $\alpha$ - najlepsza (najwyższa) ocena, jaką MAX może już na pewno zagwarantować w przeszukanej części drzewa
    \item $\beta$ - najlepsza (najniższa) ocena, jaką MIN może już na pewno zagwarantować
\end{itemize}
 
Gdy w trakcie przeszukiwania zachodzi warunek $\beta \le \alpha$, dalsze rozwijanie tej gałęzi nie ma sensu - jest oczywiste, 
że rodzic czyli przeciwnik nigdy nie wybierze tej ścieżki, ponieważ jego dotychczasowo poznana opcja jest co najmniej tak dobra 
jak obecnie analizowana. Przycinanie nie zmienia ostatecznie wybranego ruchu tylko skraca obliczenia.
 
Pełne drzewo Minimax o głębokości d=4 i rozgałęzieniu 4 (cztery możliwe kierunki ruchu) miałoby teoretycznie $4^4 = 256$ liści. 
Po zastosowaniu przycinania alfa-beta liczba realnie ocenianych węzłów spada do rzędu $\mathcal{O}(b^{d/2})$ 
(w optymalnym przypadku uporządkowania ruchów), co w praktyce pozwala na uruchamianie Minimax płynnie w czasie rzeczywistym.

\subsection{Implementacja w pliku enemy.py}

Klasa EnemyTank dziedziczy po klasie bazowej Tank i rozszerza ją o logikę decyzyjną opartą 
na Minimax.

\lstinputlisting[caption= Implementacja Minimax z przycinaniem alfa-beta (fragment z enemy.py), label=minimax]{src/enemy.py}
 
Ze względu na koszt obliczeniowy Minimax (rzędu kilkuset wywołań \_eval na jedno wywołanie \_best\_direction),
uruchomienie algorytmu w każdej klatce gry powodowałoby spadki wydajności. Zastosowano więc dwie optymalizacje czasowe:
 
\begin{itemize}
    \item \textbf{Cykl decyzyjny} - Minimax uruchamiany jest co $\text{THINK\_EVERY} = 15$ klatek, a przez pozostałe
     14 klatek bot kontynuuje ruch w ostatnio wybranym kierunku. Daje to efektywny krok decyzyjny, wystarczająco szybki 
     dla gracza.
    \item \textbf{Losowy offset} - każdy bot ma indywidualny losowy offset czasowy (random.randint(0, THINK\_EVERY-1)), 
    więc boty nie myślą wszystkie w tej samej klatce.
\end{itemize}
 
Dodatkowo wykryto, żeprzy długiej grze bot może 'utknąć' w patowych konfiguracjach (np. między dwoma ścianami z Minimax'em 
wskazującym kierunek niemożliwy do zrealizowania). Zaimplementowano detekcję stanu utknięcia: 
jeśli pozycja bota nie zmieniła się przez 30 klatek, wymuszany jest wybór losowego wolnego kierunku.
 
Po wybraniu nowego kierunku przez Minimax, zanim bot zacznie się w nim poruszać, wykonywane jest tzw. autocentrowanie do osi 
kafelka prostopadłej do kierunku ruchu: przy skręcie w poziom wyrównywana jest pozycja pionowa Y
bota do środka kafelka, przy skręcie w pion - pozycja pozioma X. Mechanizm działa z tolerancją 12 pikseli i zapobiega 
dryfowi między kafelkami, który mógłby spowodować, że bot wjedzie w konfigurację, z której nie da się ruszyć dalej w wybranym kierunku.
 
Niezależnie od powyższego, w sytuacji gdy bot mimo wszystko wjedzie w ścianę, stosowany jest drugi mechanizm awaryjny - snap do 
środka kafelka po kolizji ze ścianą.

\section{Tryb Capture Point - algorytm A*}
\subsection{Idea algorytmu A*}

A* jest algorytmem wyszukiwania optymalnej ścieżki w grafie ważonym, łączącym zalety algorytmu Dijkstry 
(gwarancja optymalności) z heurystycznym przeszukiwaniem 'najlepszy najpierw' (efektywność czasowa).
 
W zaimplementowanym wariancie graf jest dyskretną siatką kafelków 20x15 (300 węzłów), a każdy węzeł jest połączony z czterema 
sąsiadami (góra, dół, lewo, prawo). Algorytm utrzymuje dwie metryki dla każdego rozważanego węzła n:
 
\begin{itemize}
    \item $g(n)$ - rzeczywisty koszt dotarcia od~węzła startowego do n
    \item $h(n)$ - heurystyczny szacunek pozostałego kosztu od n do  celu
\end{itemize}
 
Priorytet rozwijania węzła to f(n) = g(n) + h(n). Kolejka priorytetowa zwraca zawsze 
węzeł o najniższym f, dzięki czemu algorytm szybko zbiega do celu, jednocześnie gwarantując optymalność 
ścieżki pod warunkiem, że heurystyka h jest dopuszczalna (nigdy nie przeszacowuje rzeczywistego kosztu).

\subsection{Heurystyka euklidesowa i wagi przejść kafelków}

Wykorzystaną heurystyką jest dystans euklidesowy:
 
\begin{equation}
h(a, b) = \sqrt{(a_x - b_x)^2 + (a_y - b_y)^2}
\end{equation}
 
Euklidesowa odległość po prostej między dwoma punktami nie może być większa niż rzeczywista trasa po siatce 
4-kierunkowej, ponieważ trasa po siatce zawsze realizuje przemieszczenie w osiach X i Y oddzielnie 
(sekwencja kroków w poziomie + kroków w pionie), a suma długości boków trójkąta prostokątnego jest co najmniej 
równa przeciwprostokątnej.
 
Koszt przejścia step cost zależy od typu kafelka:
 
\begin{table}[H]
    \centering
    \caption{Wagi przejść kafelków w~algorytmie A\textsuperscript{*}}
    \label{tab:astar_weights}
    \begin{tabular}{|l|c|l|}
        \hline
        \textbf{Typ kafelka} & \textbf{Koszt} & \textbf{Interpretacja} \\
        \hline
        Wolne pole              & 1   & przejście natychmiastowe \\
        Zniszczalna ściana      & 6   & wymaga rozwalenia (czas + pocisk) \\
        Niezniszczalna ściana   & $\infty$ & kafelek niedostępny \\
        \hline
    \end{tabular}
\end{table}
 
Wartość 6 dla zniszczalnej ściany jest dobrana z rozmysłem: bot wybierze obejście wolnymi polami, 
dopóki obejście wymaga mniej niż 6 dodatkowych kafelków. Gdy obejście byłoby znacznie dłuższe, opłaca się rozwalić 
ścianę. Naturalnie tworzy to zachowanie typu 'rozwalam ścianę tylko, gdy to się opłaca'.

\subsection{Cache'owanie ścieżki i warunki przeliczania}

Algorytm A* nie jest uruchamiany co klatkę - byłoby to zbyt kosztowne i niepotrzebne. 
Obliczona ścieżka jest przechowywana w polu self.path jako lista współrzędnych kafelków i przeliczana tylko wtedy, gdy:
 
\begin{itemize}
    \item Zmienił się cel (punkt przeniósł się do innego sektora po przekroczeniu progu 25\% postępu)
    \item Zmienił się układ niezniszczalnych ścian (zniszczona ściana otwiera nowe trasy)
    \item Lista ścieżki została wyczerpana (bot dotarł do kolejnego węzła końcowego)
\end{itemize}
 
Po dotarciu do każdego kafelka ścieżki ten kafelek jest usuwany z self.path, a bot kieruje się 
do następnego. W ten sposób, mimo iż A* jest kosztowny obliczeniowo, jego średnie obciążenie procesora 
jestniskie (typowo kilka wywołań na sekundę).

\subsection{Implementacja w pliku enemy2.py}

Listing poniżej przedstawia główną funkcję implementującą algorytm A*.

\lstinputlisting[caption=Implementacja algorytmu A* (fragment z enemy2.py), label=astar]{src/enemy2.py}
 
Strukturą danych przechowującą węzły do rozwinięcia jest kolejka priorytetowa zaimplementowana przez moduł heapq z biblioteki 
standardowej Pythona. Słownik came\_from zapamiętuje rodzica każdego węzła, co umożliwia rekonstrukcję ścieżki 
przez cofanie się od węzła docelowego do startowego po zakończeniu przeszukiwania.
 
Zachowanie bota wewnątrz strefy przejmowania nie wykorzystuje już A* - gdy bot dotrze do strefy, losuje losowy 
punkt patrolu wewnątrz strefy i jedzie do niego. Jeśli gracz również znajduje się w strefie, celem patrolu staje się 
sam gracz (bot przechodzi w tryb agresywny). Po dotarciu do celu losowany jest nowy punkt, aby bot nie stał w miejscu i nie 
był łatwym celem.
 
Dodatkowo, gdy bot jest jeszcze w drodze dostrefy (tryb \_drive\_to\_point), uwzględniana jest obecność power-upów 
na mapie. Jeśli najbliższy power-up znajduje się bliżej bota (w metryce Manhattan) niż środek strefy przejmowania, 
bot tymczasowo zmienia cel A* naokolicę tego power-upa (jego rect rozszerzony o 60 pikseli w każdą stronę). Pozwala 
to botowi zbierać bonusy po drodze bez rezygnacji z głównego celu, ale nie przerywa już rozpoczętego przejmowania - logika 
ta działa wyłącznie poza strefą.

\section{Heurystyki uzupełniające}
\textbf{Wykrywanie linii strzału}

Wykrywanie linii strzału (Line Of Sight, LOS) wykorzystywane jest w dwóch miejscach algorytmu 
zróżnymi stopniami restrykcyjności.
 
W funkcji oceny Minimax \_eval sprawdzenie LOS odbywa się przez rzucenie wirtualnego promienia ze stałym
krokiem STEP=20 pikseli od pozycji bota do pozycji gracza. Promień blokowany jest wyłącznie przez niezniszczalne 
ściany - zniszczalne traktowane są jak potencjalna 'zasłona dymna' (bot może je rozwalić strzałem, 
więc obecność takiej ściany nie jest sygnałem braku LOS).
 
Natomiast w rzeczywistej decyzji o oddaniu strzału \_can\_shoot,
sprawdzenie jest ściślejsze: promień rzucany jest ze mniejszym krokiem 10 pikseli i blokowany przez obie grupy ścian 
(zniszczalne i niezniszczalne). Logika ta zapobiega marnowaniu pocisków, które rozwaliłyby ścianę, ale nie trafiłyby w gracza.

\textbf{Unikanie friendly fire i marnowania pocisków}

Zaimplementowano dwa mechanizmy zapobiegające bezsensownym strzałom:

\begin{itemize}
\item Friendly fire avoidance - przed oddaniem strzału bot sprawdza, czy w jego linii ognia nie 
stoi inny bot (sojusznik). Funkcja \_ally\_in\_line sprawdza wszystkich sojuszników w zasięgu 200 pikseli na osi 
kierunku ruchu, z tolerancją prostopadłą 22 pikseli (czołg ma 20px szerokości, dodatkowa tolerancja 
zabezpiecza przed strzałem w bok sojusznika).
\item Friendly bullet detection - przed oddaniem strzału w zniszczalną ścianę bot sprawdza, czy w tym samym kierunku 
nie leci już pocisk z innego źródła (innego bota lub gracza). Funkcja \texttt\_friendly\_bullet\_in\_path szuka pocisków 
o tym samym kierunku w zasięgu 70 pikseli przed botem. Mechanizm ten zapobiega sytuacji, w której bot strzela w ścianę, którą 
za moment zniszczy już lecący pocisk - jego własny pocisk poleciałby wtedy do końca ekranu w pustą przestrzeń.
\end{itemize}

\textbf{Zachowanie kolizyjne}

System obsługi kolizji obejmuje dwa nieoczywiste mechanizmy:
 
\begin{itemize}
\item Snap do~środka kafelka - gdy bot po zniszczalnej kolizji wciąż znajduje się wewnątrz ściany 
(np. wskutek faktu, że jego hitbox po rotacji ma rozmiar 34x34 pikseli a nie 20x20, pozycja bota jest wyciągana do środka 
aktualnego kafelka, czyli wielokrotności 40 plus 20). Środek kafelka to jedyny punkt w siatce, w którym 34-pikselowy hitbox 
ma gwarancję zmieszczenia się w 40-pikselowym kafelku bez stykania się ze ścianami sąsiednich kafelków.
\item Unik od sojusznika - gdy bot wpadnie na innego bota, cofa się do pozycji sprzed ruchu i wybiera deterministycznie 
kierunek prostopadły do osi sojusznika (jeśli zderzenie nastąpiło w osi X, ucieka w osi Y, i odwrotnie). Wybór 
jest deterministyczny, a nie losowy, oraz ma cooldown 20 klatek - bez niego boty wpadające na siebie kręciły się 
wokół własnej osi co klatkę.
\end{itemize}

\section{Napotkane problemy i sposoby ich rozwiązania}

W trakcie implementacji projektu napotkano i rozwiązano szereg problemów. Najważniejsze z nich opisano poniżej.
 
\textbf{Problem 2: Boty utykały w ścianach po rotacji.} 

W pliku tank.py obraz czołgu base\_image jest tworzony 
jako powierzchnia 34x34 piksele - musi pomieścić korpus 22x22, gąsienice wystające po bokach oraz lufę wystającą ku górze. 
Jednocześnie w konstruktorze klasy Tank hitbox kolizyjny self.rect jest jawnie ustawiony na 20x20 pikseli. 
Problem pojawia się przy pierwszym wywołaniu metody rotate, która zastępuje hitbox bounding-boxem obrazu: 
self.rect = self.image.get\_rect(...). W efekcie hitbox bota po pierwszym obrocie urasta do 34x34 pikseli 
(rozmiaru obrazu) i nie wraca już do pierwotnego rozmiaru. W konsekwencji bot mógł wjechać w konfigurację, 
z której nie dało się wyjść standardowym snapem do rogu kafelka, ponieważ 34-pikselowy hitbox w rogu zawsze styka się 
ze ścianami sąsiednich kafelków. 

Rozwiązanie: snap do środka kafelka (współrzędne typu 20 + 40k), a nie do rogu - tylko środek kafelka gwarantuje, 
że 34-pikselowy hitbox zmieści się w 40-pikselowym kafelku bez stykania się ze ścianami sąsiednich kafelków (margines 
3 piksele z każdej strony).
 
\textbf{Problem 3: Boty kręciły się wokół własnej osi.} 

Bot wykonywał obrót lufy do gracza w każdej klatce, w której Minimax 
nie był uruchamiany (przez 14 klatek z 15-klatkowego cyklu), niezależnie od tego, czy reload pocisku był gotowy. 
Powodowało to migotanie wizualne i obrazanie się bota. 

Rozwiązanie: rotacja lufy wykonywana tylko, gdy get\_reload\_progress() >= 1.0 (czyli pocisk jest gotowy do strzału).
 
\textbf{Problem 5: Boty wpadające na~siebie kręciły się.} 

Gdy bot zderzył się z innym botem, każda kolejna klatka kolizji wybierała losowy nowy kierunek z listy preferowanych. 
Skutkiem był efekt nieskończonej rotacji bez rzeczywistego ruchu. 

Rozwiązanie: deterministyczny wybór kierunku (zawsze ten sam dla danego ustawienia botów) plus cooldown 20 klatek 
przed kolejną zmianą kierunku.